% chapters/requirements/rf_courses.tex
% RF-3: Cursos académicos
% ============================================================================

\item \textbf{Cursos académicos.} \label{rf:courses}
El sistema se organiza en torno a cursos académicos que pertenecen a una organización. En
todo momento existe exactamente un curso activo por organización; los cursos anteriores se
archivan y solo permanecen accesibles en modo lectura para los gestores. La transición de
un curso al siguiente es una operación semiautomática ejecutada por un gestor que archiva
el curso actual, desactiva a los usuarios no gestores y prepara el sistema para la
importación de los datos del nuevo curso.

\begin{functional}

    % RF-3.1 ----------------------------------------------------------------
    \item \textbf{Creación de curso académico.} \label{rf:courses:create}
    El sistema debe permitir a un gestor crear un nuevo curso académico para su
    organización, activándolo como curso vigente.

    \textbf{Entrada:} petición POST al recurso de cursos con la etiqueta del curso (por
    ejemplo, <<2025--2026>>), la fecha de inicio y la fecha de fin (opcional).

    \textbf{Proceso:} el sistema verifica que no exista ya un curso con la misma etiqueta
    dentro de la organización. El nuevo curso se marca como activo
    (\texttt{is\_active=True}). Si existía un curso activo anterior en la misma
    organización, se ejecuta automáticamente el proceso de transición descrito en
    \cref{rf:courses:transition}. Las fechas de inicio y fin son informativas y no disparan
    automatismos.

    \textbf{Salida:} respuesta con los datos del curso creado (etiqueta, fechas, estado
    activo) y código \acs{http}~201.

    % RF-3.2 ----------------------------------------------------------------
    \item \textbf{Modificación de curso académico.} \label{rf:courses:update}
    El sistema debe permitir a un gestor modificar los datos de un curso académico.

    \textbf{Entrada:} petición PATCH al recurso del curso con los campos a modificar
    (etiqueta, fecha de inicio, fecha de fin).

    \textbf{Proceso:} el sistema valida los datos actualizados y registra el cambio en el
    \textit{log} de auditoría con los valores anteriores y nuevos.

    \textbf{Salida:} respuesta con los datos actualizados del curso.

    % RF-3.3 ----------------------------------------------------------------
    \item \textbf{Transición de curso.} \label{rf:courses:transition}
    El sistema debe ejecutar el proceso de transición cuando se crea un nuevo curso
    académico, archivando el curso anterior.

    \textbf{Entrada:} evento interno disparado por la creación de un nuevo curso
    (\cref{rf:courses:create}).

    \textbf{Proceso:} el sistema realiza las siguientes acciones de forma transaccional:
    (1)~marca el curso anterior como archivado (\texttt{is\_active=False},
    \texttt{status=ARCHIVED}); (2)~desactiva (\texttt{is\_active=False}) a todos los
    usuarios no gestores de la organización; (3)~marca todas las
    \texttt{SubjectMembership} asociadas al curso archivado como
    \texttt{is\_active=False}; (4)~archiva las notificaciones del curso anterior;
    (5)~establece todas las instancias de examen que no estén ya finalizadas o archivadas
    al estado archivado, independientemente de su estado actual. Se registra la operación
    completa en el \textit{log} de auditoría.

    \textbf{Salida:} el curso anterior queda archivado y el sistema está preparado para
    recibir las importaciones de usuarios y asignaturas del nuevo curso.

    % RF-3.4 ----------------------------------------------------------------
    \item \textbf{Consulta de cursos.} \label{rf:courses:list}
    El sistema debe permitir al gestor consultar la lista de cursos académicos de la
    organización, diferenciando el activo de los archivados.

    \textbf{Entrada:} petición GET al recurso de cursos con filtros opcionales.

    \textbf{Proceso:} el sistema recupera los cursos de la organización que cumplen el
    filtro indicado.

    \textbf{Salida:} respuesta con la lista de cursos, incluyendo etiqueta, fechas, estado
    y número de asignaturas asociadas a cada uno.

    % RF-3.5 ----------------------------------------------------------------
    \item \textbf{Navegación entre cursos por gestor.} \label{rf:courses:navigate}
    El sistema debe permitir a un gestor consultar los datos de cualquier curso archivado
    de su organización, incluyendo sus asignaturas, exámenes e instancias.

    \textbf{Entrada:} petición GET a los recursos de asignaturas, exámenes o instancias
    incluyendo el identificador del curso como parámetro de contexto.

    \textbf{Proceso:} el sistema filtra los resultados al curso indicado, que debe
    pertenecer a la organización del gestor. Las operaciones de lectura se permiten con
    normalidad; cualquier operación de escritura sobre un curso archivado se rechaza.

    \textbf{Salida:} respuesta con los datos solicitados del curso seleccionado. Si se
    intenta una operación de escritura sobre un curso archivado, respuesta de error con
    código \acs{http}~403.

\end{functional}
